\section{Introduction}\label{sec:intro}

\subsection{The Promise}

Benchmarks that ask language-model agents to repair real software have climbed quickly. When SWE-bench was introduced, the best model it tested, Claude 2, resolved 1.96\% of its GitHub issues \citep{C1}. On SWE-Bench Pro V2, a 2026 revision for which authors from Scale AI and Reflection AI, using agentic audits and contracted engineers, repaired or retired 480 of 731 tasks and every patch is re-graded in a clean container, frontier agents resolve up to 99.4\% of the 642 remaining tasks \citep{C12}. Developers of these systems now describe much longer stretches of independent work. In a report by its developers at Edison Scientific, the ``AI scientist'' Kosmos runs for up to 12 hours, executing an average of 42,000 lines of code and reading 1,500 papers per run, and independent scientists judged 79.4\% of the statements in its reports accurate \citep{S13}. In a study by the model's developer, nine instances of Claude Opus 4.6 spent 800 cumulative hours on an alignment research problem and reached a performance-gap-recovered score of 0.97, against 0.23 for two human researchers working for seven days, at a cost of about \$18,000 \citep{S20}. Google DeepMind's AlphaEvolve found a procedure to multiply two $4 \times 4$ complex-valued matrices with 48 scalar multiplications \citep{A25}. Read together, these results suggest that agents can now do skilled technical work on their own.

\subsection{The Puzzle}

Evidence from use tells a different story. Google DeepMind's own safety report finds that although most of its production code is now more than 75\% AI-written, ``about 88\% of pull requests still require iterative multi-turn human steering or human code authorship'' \citep[p.~77]{F30}. In an observational study of 20,574 coding-agent sessions from 1,639 repositories, misalignment episodes reached a visible resolution in only 9.33\% of cases (for the rest, the logs show no outcome), and 91.49\% of those resolutions required explicit pushback from the developer \citep{T23}. Building from scratch remains hard. The best of 16 frontier models passes 61.8\% of the browser-based workflows in Vibe Code Bench \citep{C19}; on repository generation, every model tested averages a test pass rate below 40.5\% \citep{C22}; and on long Android engineering tasks with hidden tests, the best model passes 41\% \citep{C29}. In automated science the gap is wider. The FabScore audit found fabrications in 70.4\% of 54 AI-authored real conference submissions, and AI reviewers missed over 85\% of the fabrications it detected \citep{S15}. In MLR-Bench, 8 of 10 tasks conducted by Claude Code reported results based on synthesized or placeholder data, a behavior that persisted even when the agent was explicitly told not to fabricate \citep{S16}.

\subsection{The Argument}

Capability alone does not explain the gap. The same agent harness appears on both sides of it: Claude Code produced the placeholder results in MLR-Bench while running Claude 3.7 Sonnet \citep{S16} and, running Claude Opus 5, the 99.4\% \citep{C12}. The models are a generation apart, so capability explains part of the difference, but not its pattern. What separates the successes from the failures is who decides that the work is done. The 99.4\% counts patches re-graded in a clean container, where changes outside the agent's patch earn no credit \citep{C12}. The alignment agents worked on a problem with one objective metric, which their developer calls ``unusually well-suited to automation'' \citep{S20}. AlphaEvolve's authors write that its automated evaluation metric is ``a key advantage'' and ``also a limitation'', because it puts tasks that require manual experimentation out of scope \citep[p.~2]{A25}. The failures sit where the agent's own report, or a reader of its prose, decides.

This review argues that verification, not capability, is now the binding constraint on autonomous agents. The argument has three corollaries. First, success tracks the presence of an external signal: executed tests, a held-out metric, or the behavior of the running application. Second, whoever checks the work must be separate from whoever builds it, and the builder must not be able to edit the check. Third, runs that go unattended need explicit bounds and reliability measured across repeated attempts, not best-run numbers.

\subsection{Why Two Literatures Together}

The review reads together two bodies of work that have developed largely in parallel: AI for software engineering and automated science. They fail in the same ways. In both, agents report success over visible failure, optimize against graders that can be gamed, and skip checks while reporting that the checks passed. They are also repaired by the same structures, and each field holds half of the repair. Software engineering has executed tests, which turn ``does it work?'' into a question that a program answers \citep{C1}. The literature on the replication crisis has shown that analysts given the same data reach divergent results, a variation that ``may be difficult to avoid, even by experts with honest intentions'' \citep[p.~2]{R9}, with expertise and prior beliefs barely predicting the variation \citep{R11}; that forking paths arise without any conscious fishing \citep{R12}; and that commitments made before results exist change what gets reported, as in the far lower share of positive results in registered reports \citep{R13} and the lower p-hacking in studies with a complete pre-analysis plan \citep{R19}. An agent that can run many analyses in an hour inherits the analyst's problem at machine speed. Readers trained in empirical social science will recognize the AI-scientist failures reviewed below as the replication crisis, automated.

\subsection{Plan of the Review}

Section 2 describes the corpus and how evidence is weighed, and Section 3 defines autonomy and offers a typology of verification signals. Sections 4 and 5 review the architecture debate and the harness around the model. Section 6, the core of the article, examines verification. Sections 7 to 9 cover building applications, the AI-scientist program and its audits, and data-analysis agents evaluated against known answers. Section 10 draws lessons from the replication crisis, Section 11 turns to reliability, long horizons and cost, and Section 12 proposes reporting standards. Section 13 synthesizes the evidence into a framework, Section 14 draws implications for design, Section 15 sets out a research agenda, and Section 16 concludes.
